\documentclass[11pt,letterpaper]{article}

% --- Packages ---
\usepackage[utf8]{inputenc}
\usepackage[T1]{fontenc}
\usepackage{amsmath,amssymb,amsthm}
\usepackage{algorithm}
\usepackage{algpseudocode}
\usepackage{geometry}
\usepackage{hyperref}
\usepackage{booktabs}
\usepackage{enumitem}
\usepackage{graphicx}
\usepackage{bm}

\geometry{margin=1in}
\hypersetup{colorlinks=true, linkcolor=blue, citecolor=blue, urlcolor=blue}

% --- Theorem environments ---
\newtheorem{definition}{Definition}[section]
\newtheorem{theorem}{Theorem}[section]
\newtheorem{proposition}{Proposition}[section]
\newtheorem{remark}{Remark}[section]

% --- Notation macros ---
\newcommand{\R}{\mathbb{R}}
\newcommand{\frame}{\mathbf{F}}
\newcommand{\pos}{\mathbf{p}}
\newcommand{\basis}{\mathbf{e}}
\newcommand{\grad}{\nabla}
\newcommand{\manifold}{\mathcal{M}}
\newcommand{\tangent}{T}
\newcommand{\embed}{\Phi}
\newcommand{\obj}{\mathcal{L}}
\newcommand{\cov}{\mathbf{C}}
\newcommand{\eye}{\mathbf{I}}
\newcommand{\norm}[1]{\left\lVert #1 \right\rVert}

\title{%
  \textbf{Intelligent Navigation of $N$-Dimensional Manifolds} \\[6pt]
  \large A Formal Specification for Manifold-Aware Descent \\
  via Turtle Geometry and Local Principal Component Analysis
}

\author{
  Eric G. Suchanek, PhD \\
  \texttt{suchanek@mac.com} \\[4pt]
  proteusPy Project \\
  \url{https://github.com/suchanek/proteusPy}
}

\date{\today}

\begin{document}
\maketitle

\begin{abstract}
We present a formal specification for navigating $N$-dimensional embedding
spaces by combining turtle geometry with local principal component analysis.
The system generalizes the classical 3D turtle to $N$ dimensions via Givens
rotations over an orthonormal frame, then couples this with adaptive,
position-dependent PCA to discover the tangent space of the data manifold
at each step. Gradients are projected onto this tangent space and weighted
by eigenvalue magnitude, yielding a form of empirical natural gradient
descent that respects the intrinsic geometry of the embedding rather than
its ambient coordinate system. We provide complete mathematical definitions,
algorithmic specifications, correctness properties, and complexity analysis.
\end{abstract}

\tableofcontents
\newpage

%======================================================================
\section{Introduction}
\label{sec:intro}
%======================================================================

Large language models and other deep learning systems embed discrete
symbolic structures into continuous vector spaces of high dimension $N$.
These embedding spaces are not isotropic: the data concentrates on
lower-dimensional manifolds $\manifold \subset \R^N$ whose intrinsic
dimensionality $d \ll N$ varies with position. Standard optimization
methods---gradient descent, Adam, and their variants---treat all $N$
dimensions equivalently, wasting capacity on off-manifold noise and
ignoring the curvature of the data manifold.

We propose an alternative: \emph{manifold-aware navigation} using a
generalized turtle geometry. The turtle carries an orthonormal frame
that is continuously re-aligned to the local tangent space of $\manifold$
via PCA on a neighborhood of data points. This transforms the optimization
problem from descent in the ambient $\R^N$ to descent along the manifold,
suppressing off-manifold components and weighting on-manifold directions
by their data variance.

This document specifies two components:
\begin{enumerate}[label=(\roman*)]
  \item \textbf{TurtleND}: An $N$-dimensional turtle with position, orthonormal frame, and Givens rotation primitives.
  \item \textbf{ManifoldWalker}: A navigation system that wraps TurtleND with local PCA steering, eigenvalue-weighted gradient projection, and adaptive dimensionality estimation.
\end{enumerate}

%======================================================================
\section{Mathematical Preliminaries}
\label{sec:prelim}
%======================================================================

\begin{definition}[Orthonormal Frame]
\label{def:frame}
An \emph{orthonormal frame} in $\R^N$ is an ordered set of $N$ vectors
$\frame = (\basis_0, \basis_1, \ldots, \basis_{N-1})$ satisfying:
\begin{equation}
  \langle \basis_i, \basis_j \rangle = \delta_{ij}
  \qquad \forall\; i, j \in \{0, \ldots, N-1\}
\end{equation}
where $\delta_{ij}$ is the Kronecker delta. Equivalently, the matrix
$F \in \R^{N \times N}$ with rows $\basis_i$ satisfies $F F^\top = \eye_N$.
\end{definition}

\begin{definition}[Givens Rotation]
\label{def:givens}
A \emph{Givens rotation} $G(i, j, \theta)$ in $\R^N$ acts on the plane
spanned by coordinates $i$ and $j$, rotating by angle $\theta$:
\begin{equation}
  G(i, j, \theta) = \eye_N + (\cos\theta - 1)(\mathbf{e}_i\mathbf{e}_i^\top + \mathbf{e}_j\mathbf{e}_j^\top) + \sin\theta(\mathbf{e}_i\mathbf{e}_j^\top - \mathbf{e}_j\mathbf{e}_i^\top)
\end{equation}
where $\mathbf{e}_i, \mathbf{e}_j$ are standard basis vectors. A Givens
rotation preserves orthonormality: if $F$ is orthonormal, so is $G \cdot F$.
\end{definition}

\begin{definition}[Data Manifold]
\label{def:manifold}
Given a set of embedding vectors $\mathcal{X} = \{\mathbf{x}_1, \ldots, \mathbf{x}_M\} \subset \R^N$,
the \emph{data manifold} $\manifold$ is the smooth submanifold (or the
best smooth approximation thereof) on which the data concentrates. The
\emph{intrinsic dimensionality} $d(\pos)$ at a point $\pos$ is the
dimension of the tangent space $\tangent_\pos\manifold$.
\end{definition}

\begin{definition}[Local Tangent Space via PCA]
\label{def:tangent}
Given a point $\pos \in \R^N$ and its $k$-nearest neighborhood
$\mathcal{N}_k(\pos) \subset \mathcal{X}$, the \emph{empirical local
tangent space} is the span of the top-$d$ eigenvectors of the local
covariance matrix:
\begin{equation}
  \cov(\pos) = \frac{1}{k-1} \sum_{\mathbf{x} \in \mathcal{N}_k(\pos)} (\mathbf{x} - \bar{\mathbf{x}})(\mathbf{x} - \bar{\mathbf{x}})^\top
\end{equation}
where $\bar{\mathbf{x}}$ is the neighborhood centroid. The eigendecomposition
$\cov = V \Lambda V^\top$ with $\lambda_1 \geq \lambda_2 \geq \cdots \geq \lambda_N \geq 0$
yields the principal directions $\mathbf{v}_1, \ldots, \mathbf{v}_N$ as
columns of $V$.
\end{definition}

%======================================================================
\section{TurtleND Specification}
\label{sec:turtlend}
%======================================================================

\subsection{State}

The TurtleND maintains:
\begin{itemize}[nosep]
  \item Position: $\pos \in \R^N$
  \item Frame: $\frame = (\basis_0, \basis_1, \ldots, \basis_{N-1})$, an orthonormal frame
\end{itemize}

By convention:
\begin{itemize}[nosep]
  \item $\basis_0$: \emph{heading} (direction of movement)
  \item $\basis_1$: \emph{left}
  \item $\basis_2$: \emph{up} (when $N \geq 3$)
  \item $\basis_3, \ldots, \basis_{N-1}$: higher-dimensional basis vectors
\end{itemize}

\noindent\textbf{Initial state:} $\pos = \mathbf{0}$, $\frame = \eye_N$.

\subsection{Operations}

\begin{table}[h]
\centering
\caption{TurtleND operations and their mathematical definitions.}
\label{tab:ops}
\begin{tabular}{@{}lll@{}}
\toprule
\textbf{Operation} & \textbf{Definition} & \textbf{Plane} \\
\midrule
$\textsc{Move}(d)$ & $\pos \leftarrow \pos + d \cdot \basis_0$ & --- \\
$\textsc{Rotate}(\theta, i, j)$ & Apply $G(i, j, \theta)$ to $\frame$ & $(\basis_i, \basis_j)$ \\
$\textsc{Turn}(\theta)$ & $\textsc{Rotate}(\theta, 0, 1)$ & heading--left \\
$\textsc{Roll}(\theta)$ & $\textsc{Rotate}(\theta, 1, 2)$ & left--up \\
$\textsc{Pitch}(\theta)$ & $\textsc{Rotate}(-\theta, 0, 2)$ & heading--up \\
$\textsc{Yaw}(\theta)$ & $\textsc{Rotate}(180° - \theta, 0, 1)$ & heading--left \\
$\textsc{ToLocal}(\mathbf{g})$ & $\ell_i = \langle \basis_i, \mathbf{g} - \pos \rangle$ & --- \\
$\textsc{ToGlobal}(\bm{\ell})$ & $\mathbf{g} = \pos + \sum_i \ell_i \basis_i$ & --- \\
\bottomrule
\end{tabular}
\end{table}

\begin{proposition}[Frame Invariant]
\label{prop:frame}
All TurtleND rotation operations preserve the orthonormality of $\frame$.
That is, if $\frame$ is orthonormal before a rotation, it is orthonormal
after (up to floating-point precision).
\end{proposition}

\begin{proof}
Each rotation modifies exactly two rows of $F$ via a $2 \times 2$ rotation
matrix $\begin{psmallmatrix} \cos\theta & \sin\theta \\ -\sin\theta & \cos\theta \end{psmallmatrix}$
applied to rows $i$ and $j$. Since the original rows are orthonormal and
a rotation matrix preserves inner products, the result is orthonormal. \qed
\end{proof}

\begin{proposition}[Coordinate Transform Invertibility]
\label{prop:coords}
$\textsc{ToGlobal}(\textsc{ToLocal}(\mathbf{g})) = \mathbf{g}$ for all $\mathbf{g} \in \R^N$.
\end{proposition}

\begin{proof}
Let $\bm{\delta} = \mathbf{g} - \pos$. Then $\textsc{ToLocal}(\mathbf{g}) = F\bm{\delta}$
and $\textsc{ToGlobal}(F\bm{\delta}) = \pos + F^\top F \bm{\delta} = \pos + \bm{\delta} = \mathbf{g}$
since $F^\top F = \eye_N$. \qed
\end{proof}

\subsection{3D Backward Compatibility}

\begin{theorem}[Equivalence]
\label{thm:compat}
TurtleND with $N=3$ produces identical position and frame vectors as
Turtle3D for all sequences of \textsc{Move}, \textsc{Turn}, \textsc{Roll},
\textsc{Pitch}, and \textsc{Yaw} operations with identical parameters.
\end{theorem}

This is verified empirically by exhaustive testing of individual operations
and compound sequences.

%======================================================================
\section{ManifoldWalker Specification}
\label{sec:walker}
%======================================================================

\subsection{State}

The ManifoldWalker maintains:
\begin{itemize}[nosep]
  \item A TurtleND instance (position and frame)
  \item An embedding matrix $\mathcal{X} \in \R^{M \times N}$
  \item An objective function $\obj : \R^N \to \R$ to minimize
  \item Parameters: neighborhood size $k$, variance threshold $\tau$, learning rate $\eta$
  \item Diagnostics: eigenvalue spectrum $\bm{\lambda}$, intrinsic dimensionality $d$
\end{itemize}

\subsection{Core Algorithm: Manifold-Aware Descent}

\begin{algorithm}[H]
\caption{ManifoldWalker.Step}
\label{alg:step}
\begin{algorithmic}[1]
\Require Current position $\pos$, embeddings $\mathcal{X}$, objective $\obj$, parameters $k, \tau, \eta$
\Ensure Updated position $\pos'$, objective value $\obj(\pos')$
\Statex

\State \textbf{Phase 1: Discover local manifold geometry}
\State $\mathcal{N} \leftarrow k\text{-nearest neighbors of } \pos \text{ in } \mathcal{X}$
\State $\bar{\mathbf{x}} \leftarrow \frac{1}{k}\sum_{\mathbf{x} \in \mathcal{N}} \mathbf{x}$
\State $\cov \leftarrow \frac{1}{k-1}\sum_{\mathbf{x} \in \mathcal{N}} (\mathbf{x} - \bar{\mathbf{x}})(\mathbf{x} - \bar{\mathbf{x}})^\top$
\State $(\lambda_1, \mathbf{v}_1), \ldots, (\lambda_N, \mathbf{v}_N) \leftarrow \text{eigendecomposition of } \cov$ \Comment{$\lambda_1 \geq \cdots \geq \lambda_N$}

\Statex
\State \textbf{Phase 2: Estimate intrinsic dimensionality}
\State $d \leftarrow \min\left\{j : \frac{\sum_{i=1}^{j} \lambda_i}{\sum_{i=1}^{N} \lambda_i} \geq \tau\right\}$

\Statex
\State \textbf{Phase 3: Orient turtle to manifold tangent space}
\State $\basis_i \leftarrow \mathbf{v}_i$ for $i = 0, \ldots, N-1$ \Comment{Align frame to principal directions}

\Statex
\State \textbf{Phase 4: Compute and project gradient}
\State $\mathbf{g} \leftarrow \grad\obj(\pos)$ \Comment{Analytic or numerical}
\State $\bm{\ell} \leftarrow F \cdot \mathbf{g}$ \Comment{Project into local frame}

\Statex
\State \textbf{Phase 5: Suppress off-manifold components and weight}
\For{$i = 0, \ldots, N-1$}
  \If{$i \geq d$}
    \State $\ell_i \leftarrow 0$ \Comment{Off-manifold: suppress}
  \Else
    \State $\ell_i \leftarrow \ell_i \cdot \frac{\lambda_i}{\lambda_1}$ \Comment{Eigenvalue weighting}
  \EndIf
\EndFor

\Statex
\State \textbf{Phase 6: Step in global coordinates}
\State $\mathbf{s} \leftarrow F^\top \cdot \bm{\ell}$ \Comment{Map weighted gradient back to global}
\State $\pos' \leftarrow \pos - \eta \cdot \mathbf{s}$ \Comment{Descend}

\Statex
\State \Return $\pos', \obj(\pos')$
\end{algorithmic}
\end{algorithm}

\subsection{Intrinsic Dimensionality Estimation}

The intrinsic dimensionality $d(\pos)$ is estimated at each step as the
number of eigenvalues needed to capture fraction $\tau$ of the total
variance:
\begin{equation}
  d(\pos) = \min \left\{ j \in \{1, \ldots, N\} \;:\; \frac{\sum_{i=1}^{j} \lambda_i(\pos)}{\sum_{i=1}^{N} \lambda_i(\pos)} \geq \tau \right\}
\end{equation}

This is \emph{adaptive}: $d$ may vary from point to point on the manifold.
Regions of high curvature or branching may exhibit higher $d$; flat
regions may have very low $d$.

\subsection{Eigenvalue Weighting as Natural Gradient}

The weighting $w_i = \lambda_i / \lambda_1$ applied to each on-manifold
gradient component implements a form of natural gradient descent. The
local covariance $\cov(\pos)$ serves as an empirical estimate of the
metric tensor of the data manifold. Directions of high data variance
(large $\lambda_i$) receive proportionally larger step sizes, while
directions of low variance (small $\lambda_i$) are damped.

This is analogous to the Fisher Information Matrix (FIM) in classical
natural gradient methods, but estimated nonparametrically from the
data geometry rather than from the model's parameter space.

\begin{remark}[Relationship to Riemannian Optimization]
The ManifoldWalker implements a discrete approximation to Riemannian
gradient descent on the data manifold $\manifold$. At each point:
\begin{enumerate}[nosep]
  \item The local PCA approximates the tangent space $\tangent_\pos\manifold$
  \item The eigenvalue-weighted projection approximates the Riemannian gradient $\grad_{\manifold}\obj$
  \item The linear step approximates the exponential map $\exp_\pos(-\eta \grad_\manifold\obj)$
\end{enumerate}
The approximation quality depends on the neighborhood size $k$ and the
local curvature of $\manifold$.
\end{remark}

%======================================================================
\section{Coordinate Transform Geometry}
\label{sec:transforms}
%======================================================================

The turtle's frame provides a moving coordinate system. At any point $\pos$
with frame $F$:

\begin{align}
  \textsc{ToLocal}(\mathbf{g}) &= F \cdot (\mathbf{g} - \pos) \label{eq:tolocal} \\
  \textsc{ToGlobal}(\bm{\ell}) &= \pos + F^\top \cdot \bm{\ell} \label{eq:toglobal}
\end{align}

In the ManifoldWalker context, after orientation:
\begin{itemize}[nosep]
  \item $\ell_0$: component along direction of maximum local variance (heading)
  \item $\ell_1$: component along second principal direction (left)
  \item $\ell_i$ for $i < d$: on-manifold components
  \item $\ell_i$ for $i \geq d$: off-manifold components (suppressed)
\end{itemize}

%======================================================================
\section{Probing the Loss Landscape}
\label{sec:probe}
%======================================================================

The \textsc{Probe} operation evaluates the objective along a single
principal direction without moving the turtle:
\begin{equation}
  \textsc{Probe}(i, \{s_1, \ldots, s_m\}) = \left\{ \obj(\pos + s_j \cdot \basis_i) \right\}_{j=1}^{m}
\end{equation}

This enables:
\begin{itemize}[nosep]
  \item Visualization of the loss landscape along individual principal components
  \item Detection of non-convexity or saddle points in specific directions
  \item Adaptive step size selection based on local curvature estimates
\end{itemize}

%======================================================================
\section{Complexity Analysis}
\label{sec:complexity}
%======================================================================

\begin{table}[h]
\centering
\caption{Computational complexity per step.}
\label{tab:complexity}
\begin{tabular}{@{}lll@{}}
\toprule
\textbf{Operation} & \textbf{Complexity} & \textbf{Notes} \\
\midrule
$k$-NN search & $O(MN)$ & Brute force; $O(N \log M)$ with KD-tree \\
Covariance matrix & $O(k N^2)$ & \\
Eigendecomposition & $O(N^3)$ & Full spectrum; $O(N^2 d)$ for top-$d$ only \\
Gradient (numerical) & $O(2N \cdot C_\obj)$ & $C_\obj$ = cost of one $\obj$ evaluation \\
Projection \& step & $O(N^2)$ & Matrix-vector products \\
\midrule
\textbf{Total per step} & $O(MN + N^3 + NC_\obj)$ & Dominated by KNN or eigendecomp \\
\bottomrule
\end{tabular}
\end{table}

For practical embedding spaces ($N \sim 768$--$4096$, $M \sim 10^4$--$10^6$),
the KNN search dominates. This can be mitigated with approximate nearest
neighbor structures (FAISS, Annoy, ScaNN). The eigendecomposition can use
truncated SVD when only the top-$d$ components are needed.

%======================================================================
\section{Correctness Properties}
\label{sec:correctness}
%======================================================================

\begin{enumerate}[label=\textbf{P\arabic*.}, ref=P\arabic*]
  \item \label{prop:ortho} \textbf{Frame Orthonormality.} The turtle's frame remains orthonormal after any sequence of rotations: $FF^\top = \eye_N$. A periodic Gram-Schmidt re-orthonormalization corrects floating-point drift.

  \item \label{prop:invert} \textbf{Coordinate Invertibility.} $\textsc{ToGlobal} \circ \textsc{ToLocal} = \mathrm{id}$ (Proposition~\ref{prop:coords}).

  \item \label{prop:compat} \textbf{3D Compatibility.} TurtleND$(3)$ is operationally identical to Turtle3D (Theorem~\ref{thm:compat}).

  \item \label{prop:descent} \textbf{Descent Property.} For a smooth, locally convex objective $\obj$ with Lipschitz gradient, and sufficiently small $\eta$, each step reduces the objective: $\obj(\pos') \leq \obj(\pos) - c \norm{\grad_\manifold \obj(\pos)}^2$ for some $c > 0$ depending on $\eta$ and the eigenvalue spectrum.

  \item \label{prop:manifold} \textbf{Manifold Alignment.} The frame's first $d$ basis vectors span the same subspace as the top-$d$ eigenvectors of the local covariance, ensuring the step direction lies in (or near) $\tangent_\pos \manifold$.
\end{enumerate}

%======================================================================
\section{Parameters and Tuning}
\label{sec:params}
%======================================================================

\begin{table}[h]
\centering
\caption{ManifoldWalker parameters and guidance.}
\label{tab:params}
\begin{tabular}{@{}llll@{}}
\toprule
\textbf{Parameter} & \textbf{Symbol} & \textbf{Default} & \textbf{Guidance} \\
\midrule
Neighborhood size & $k$ & 50 & $\sim 5$--$10\times$ expected $d$ \\
Variance threshold & $\tau$ & 0.95 & Lower $\to$ more aggressive dimensionality reduction \\
Learning rate & $\eta$ & 0.01 & Scale inversely with gradient magnitude \\
Gradient epsilon & $\varepsilon$ & $10^{-5}$ & For numerical differentiation only \\
\bottomrule
\end{tabular}
\end{table}

%======================================================================
\section{Discussion}
\label{sec:discussion}
%======================================================================

\subsection{Why Not Cosine Distance?}

Cosine distance treats the embedding space as isotropic---all angular
directions are equally meaningful. This is false for learned embeddings.
The variance along different principal directions can differ by orders of
magnitude, and those directions \emph{rotate} as one moves through the space.
The ManifoldWalker's local PCA captures this anisotropy directly.

\subsection{Relationship to Existing Methods}

\begin{itemize}[nosep]
  \item \textbf{Natural Gradient Descent} (Amari, 1998): Uses the Fisher
    information matrix; ManifoldWalker uses the empirical data covariance
    as a nonparametric analog.
  \item \textbf{Local Tangent Space Alignment} (Zhang \& Zha, 2004):
    Discovers manifold structure via local PCA; ManifoldWalker adds
    navigation and optimization.
  \item \textbf{Riemannian Optimization} (Absil et al., 2008): Requires
    knowing the manifold geometry a priori; ManifoldWalker discovers it
    empirically at each step.
\end{itemize}

\subsection{From Molecular Geometry to Embedding Spaces}

The TurtleND originates from Turtle3D, a tool for building molecular
structures via local coordinate systems (Suchanek, 1990). The generalization
to $N$ dimensions and coupling with local PCA creates a bridge between
geometric molecular modeling and high-dimensional machine learning---the
same frame-based navigation that builds proteins can now navigate the
latent spaces that represent them.

%======================================================================
\section{Conclusion}
\label{sec:conclusion}
%======================================================================

The TurtleND + ManifoldWalker system provides a geometrically principled
framework for navigating high-dimensional embedding spaces. By continuously
re-aligning a carried orthonormal frame to the local data manifold via PCA,
it achieves manifold-aware descent that:
\begin{enumerate}[nosep]
  \item Adapts to locally varying intrinsic dimensionality
  \item Suppresses off-manifold gradient noise
  \item Weights on-manifold directions by data variance (natural gradient)
  \item Provides interpretable diagnostics (eigenvalue spectra, intrinsic dim)
  \item Enables directional probing of the loss landscape
\end{enumerate}

This replaces the assumption of isotropic space with empirical geometry,
offering a new paradigm for optimization in learned embedding spaces.

\end{document}
